\section{Open problems}\label{gl:sec:op}



\begin{problem}[The Thorin property without the zeros]\label{gl:op:direct}
Establish condition (i) of Theorem~\ref{gl:thm:family} for $\mathbb Q$, or for cyclotomic fields, without first knowing (ii). In the master theorem route, settle the parabolic distortion of Section~\ref{gl:sec:rmt-open}, uniformly in the field, and find a substitute for the residue bookkeeping at the zeros. \emph{Status.} Open. For $\mathbb Q$ at $u_0=0$ the statement to be proved is RH, so a growth bound that is itself proved under RH, such as Lemma~\ref{gl:lem:stirling}, cannot be an input to its proof.
\end{problem}

\begin{problem}[Seven with real zeros]\label{gl:op:exceptional}
Show that $K=7$ follows from the Thorin property at tilt $u_0=1$ for cyclotomic fields, that is from GRH with possible real zeros. Section~\ref{gl:sec:K} locates the arcs where a real zero can exceed the bound $N^{3/4}$; what remains is to move them into the major arcs and to redo the main term with several real zeros, checking that the constant $2.7895$ survives. In the paper of Heath-Brown and Puchta~\cite{HBP02} the major arc estimate that gives this constant (their Lemma~7) is stated for $K\ge9$, with zero density estimates in place of GRH, and the paper does not say how it is used for $K=7$; the computation of the constant itself holds for $K\ge7$. \emph{Status.} Open. Only the location of the arcs is done, and it is a heuristic estimate.
\end{problem}

\begin{problem}[The moment identity at scale]\label{gl:op:numerics}
Anchoring the tail to the exact zero count, equation~\eqref{gl:eq:anchored}, already brings the check of $p_1$ from $2.6\times10^{-6}$ to $6.3\times10^{-9}$ with $799$ zeros. Carry it to large height with Odlyzko's tables~\cite{OdlyzkoTables}, where the remaining term of~\eqref{gl:eq:anchored}, the integral of $S(t)\,t^{-2j-1}$ over the tail, sets the limit and calls for an explicit formula for $\int S$, and to $L(s,\chi)$ with the zeros of Platt~\cite{Platt16}. \emph{Status.} Partly done. The script \texttt{gl\_zeta\_moments.py} (\texttt{code/\allowbreak README.md}) reproduces the anchoring with the first $1700$ ordinates at the cuts $N=200$ and $N=799$ and at every cut $N=50,\dots,1699$, at $100$ digits (\texttt{data/gl\_zeta\_moments.json}). For $p_1$ the anchored tail typically gains about a third of the height of the cut over the smooth tail (the first and ninth deciles of the gain are $42$ and $2700$), but it gains nothing where $|S|$ is small at the cut: the gain is below a factor $10$ at $3.2\%$ of the cuts. Large heights and the characters are open.
\end{problem}

% CHECK: the model and the errors are those of code/gl_density_pf.py (full mode, data/gl_density_pf.json, official). The values 3.7575 and 3.7097 are the function side of data/gl_dirichlet.json (3.7574676509 and 3.7097071713; the zero side differs by 1.3e-5 and 3.6e-5, the tail), confirmed in misc/reviews/nick-linnik/checks-L2.md.
\begin{problem}[The negative moment as a family statistic]\label{gl:op:negmoment}
The value $\E[T^{-1/2}]=6.7333$ (and $3.7575$, $3.7097$ for $T_{\chi_4}$, $T_{\chi_5}$) is computed from $\xi$ alone: $\E[T^{-1/2}]=\Gamma(\tfrac12)^{-1}\int_0^\infty u^{-1/2}\prod_\gamma(1+u/\gamma^2)^{-1}\,du$, the product being over the ordinates $\gamma>0$. It is mostly a smooth statistic. Let $E_m$ be this integral with the first $m$ ordinates exact and the factors of the others replaced by the smooth model $\exp\big(-\tfrac1\pi\int_{t_m}^\infty\log(1+u/t^2)\,\vartheta'(t)\,dt-a_m\log(1+u/t_m^2)\big)$, where $\vartheta$ is the Riemann Siegel theta function, $t_m$ is the midpoint of $\gamma_m$ and $\gamma_{m+1}$, and $a_m=N_0(t_m)+1-m$, with $N_0=\vartheta/\pi$, is the mass of the atom at $t_m$ that the anchored tail~\eqref{gl:eq:anchored} adds. For $m=0$ the cut is $t_0=9.6669$, the point where $N_0+1=0$, and there is no atom. The relative error $E_m/\E[T^{-1/2}]-1$ is $-7.0\times10^{-3}$ for $m=0$, $-2.5\times10^{-3}$ for $m=1$, $1.0\times10^{-4}$ for $m=2$, $-3.9\times10^{-4}$ for $m=5$ and $1.4\times10^{-5}$ for $m=20$. The atom matters: cutting at the midpoint without it gives $-5.5\times10^{-3}$, $-1.2\times10^{-3}$, $-1.4\times10^{-3}$ and $-6.5\times10^{-4}$ for $m=1$, $2$, $5$, $20$, and cutting at the point where $N_0+1=m$ without it gives $-2.5\times10^{-3}$, $9.1\times10^{-5}$, $-4.0\times10^{-4}$ and $1.0\times10^{-5}$. The elasticity $(\gamma_k/E)\,\partial E/\partial\gamma_k$ of $E=\E[T^{-1/2}]$ with respect to one ordinate, the others fixed in the integral above, is $0.1799$ for $\gamma_1$ and falls to $0.0319$ for $\gamma_6$ (over all ordinates the elasticities sum to $1$, by scaling). A single function therefore says little, and the natural comparison with random matrices is in families. Compute the average of $\E[T_\chi^{-1/2}]$ over real primitive characters of conductor in $[Q,2Q]$ and compare it with the prediction of the symplectic ensemble. Since $\E[T^{-1/2}]$ is not a linear statistic of the zeros, the one level density of Katz and Sarnak~\cite{KatzSarnak99} alone does not give this prediction, and all the correlation functions enter. The circular unitary comparisons of Section~\ref{ag:sec:margins} and the moments of Keating and Snaith~\cite{KS00} (Section~\ref{ag:sec:ks}) are the models for a single function. \emph{Status.} Partly done: the single function is computed, with the model above, by \texttt{gl\_density\_pf.py}; the family average is open.
\end{problem}

\begin{problem}[The family mixture]\label{gl:op:family}
Let $T^{(q)}$ be the variable of a uniformly chosen member of $\mathcal F_q$ (Section~\ref{gl:sec:family}, under GRH for the members and with their central values nonzero). Is $T^{(q)}$ infinitely divisible, and does any $q$ with two or more members avoid the odd configuration of Proposition~\ref{gl:prop:mixture}(b), which is stated for two clocks, for every pair of its members? \emph{Status.} Open. That $T^{(q)}$ is not a GGC is settled for $q=5$ and every weight, under GRH for the two members (their values at $\tfrac12$ are nonzero numerically; Proposition~\ref{gl:prop:mod5}), and holds, under the same hypotheses, for every $q$ whose family has exactly two members, when these are in the odd configuration (Proposition~\ref{gl:prop:mixture}(b)). It does not decide whether $T^{(5)}$ is infinitely divisible.
\end{problem}

\begin{problem}[Sign definite explicit formulae]\label{gl:op:signdefinite}
Under the hypothesis of Corollary~\ref{gl:cor:K7} the explicit formula for $\psi(x;q,1)$, which up to the factor $\varphi(q)$ counts the prime ideals of $\mathbb Q(\zeta_q)$, becomes an integral against the nonnegative measure $U_{F_q,0}$, and Theorem~\ref{gl:thm:sinsq} supplies the positive family $\pi^{-1}\Re(\xi_{F_q}'/\xi_{F_q})(\alpha+ib)$ joining it to the primes. For $a\not\equiv1\pmod q$ the characters enter the formula for $\psi(x;q,a)$ with the weights $\bar\chi(a)$, which are complex or negative, and no sign is available. These are zero counting measures and their Poisson smoothings, so any gain must come from structure beyond counting, such as monotonicity or log concavity of the associated GGC densities~\cite[Ch.~4]{Bondesson}, transported to bounds on $L'/L$ near the line. Is there a sign definite explicit formula for $\psi(x;q,a)$ with $a\not\equiv1\pmod q$, in terms of such measures? \emph{Status.} Open.
\end{problem}

\begin{problem}[Certified backward continuation]\label{gl:op:backward}
The Thorin density $\pi^{-1}\Re(f'/f)(\alpha+ib)$ of Theorem~\ref{gl:thm:sinsq}, which is $\rho_{\alpha-1/2}(b)$ for $f=\xi$, is computable from the primes, with rigorous error bounds, for $\alpha>1$, and its positivity on every line $\Re s=\alpha'$ with $\alpha'>\alpha$ is a zero free half plane $\Re s>\alpha$. The book certifies the minimum at the origin for $\xi$, $\chi_4$ and $\chi_{12}$ for $\alpha<1$ up to a large height, from the zeros (Theorem~\ref{thm:dbeta:minall}). How far into the strip can positivity be certified from the primes alone, using smoothed prime sums of the kind that located the zeros of $L(s,\chi_4)$ in Section~\ref{gl:sec:sinsq}, before the ill posedness of the backward step takes over? Chapter~\ref{chap:tf} measures the cost of the step: by Theorem~\ref{tf:thm:deconv} it is a deconvolution, which is unbounded for the basepoint family, and Proposition~\ref{tf:prop:mellin} counts the primes it needs, up to about $\eta^{-1/(\alpha-\sigma_{\max})}$ for accuracy $\eta$; for the heat flow the corresponding step is stable but cannot carry real zeros (Section~\ref{tf:sec:heat}). Section~\ref{dc:sec:weil} of Chapter~\ref{chap:dc} reads the same step as a deconvolution of the explicit formula. \emph{Status.} Partly done. For $\chi_4$ the prime side~\eqref{gl:eq:Malpha} of $\pi$ times the density at $b=0$, summed over the prime powers up to $2\times10^6$ without smoothing, differs from $(\xi'/\xi)(\alpha,\chi_4)$ by $1.3\times10^{-4}$, $4.8\times10^{-3}$, $4.1\times10^{-2}$ and $8.3\times10^{-2}$ at $\alpha=1$, $\tfrac34$, $\tfrac35$ and $\tfrac{11}{20}$, that is by $0.2\%$, $12\%$, $260\%$ and $1100\%$ of the value, so that for $\alpha\le1$ the minima computed in Section~\ref{gl:sec:sinsq} come from the zeros and not from the primes. What smoothed prime sums can certify below $\alpha=1$ is open.
\end{problem}

\begin{problem}[The arithmetic L\'evy measure and the minor arcs]\label{gl:op:arith}
For $F=\mathbb Q$ the atoms of $\nu_F$ in~\eqref{gl:eq:nuF} at the primes $p$ have mass $\log p$, the weights that enter $\mathcal S(\alpha)$; the prime powers of~\eqref{gl:eq:nuF} do not enter $\mathcal S(\alpha)$. Can $\mathcal S(\alpha)$ on the minor arcs be written as an expectation against the GGC laws above? Such a representation would be a different route to the sup bound in step~(3); nothing in the chapter shows that it exists. \emph{Status.} Open.
\end{problem}

\begin{problem}[Zero free half planes in the circle method]\label{gl:op:quasi}
Assume $L(s,\chi)\ne0$ for $\Re s>\theta_0$ and every Dirichlet character. What is the best minor arc exponent $\theta(\theta_0)$ obtainable by combining the explicit formula with Vaughan's identity? The one line balance of Section~\ref{gl:sec:oct} gives $\tfrac12(1+\theta_0)$, which is better than the unconditional exponent $\tfrac{263}{308}$ of step~(3) exactly when $\theta_0<\tfrac{109}{154}$, so it does nothing at $\theta_0=\tfrac78$; it is not shown to be optimal. Does a hybrid bound, using the half plane only for the characters of small conductor, improve $\lambda$ in step~(4)? And does the half plane, which excludes exceptional zeros of large conductor, raise the constant $2.7895$ of step~(1)? The paper of Heath-Brown and Puchta~\cite{HBP02} suggests that it does not by itself: they state that zeros close to $\sigma=1$ play no role and that no explicit numerical zero free region is needed, the constant is a sum over the divisors $d\le5$ (the improvement of Elsholtz to $2.96169$ that they report takes $d\le21$ and a better sieve constant), and the restriction $q\le P=N^{45/154-4\varpi}$ of their major arcs comes from zero density estimates, whose critical zeros have $\sigma\le\tfrac34$ and are not removed by a half plane at $\tfrac78$. That restriction is not sharp for those estimates: the exponent $(\tfrac{45}{77}\kappa(\sigma)-1)(1-\sigma)$ of their Section~3 has its maximum $0.1050$ at $\sigma=0.676$, against the bound $\tfrac{31}{154}=0.2013$ that they use. \emph{Status.} Open.
\end{problem}

\begin{problem}[Finite order total positivity]\label{gl:op:pf}
RH holds if and only if the density of $Y=\sqrt{2T}\,Z$ is a P\'olya frequency function of infinite order (Proposition~\ref{pol:prop:pf}, where the variable is called $X$; Section~\ref{gl:sec:rh}). For this one density the question whether a finite order already implies RH has a trivial answer. The orders decrease, so if RH holds every such implication is true, and if RH fails the property fails at some finite order, where the implication is vacuous. The question has content for a class of functions. Let $\mathcal E$ be the class of even entire functions $F$ with $F(0)=1$ and the order, type and reflection symmetry of $\xi(\tfrac12+s)/\xi(\tfrac12)$ for which $1/F$ is the characteristic function of a probability density $f_F$; the density of $Y$ is $f_F$ for $F=\xi(\tfrac12+\cdot)/\xi(\tfrac12)$ (Proposition~\ref{pol:prop:pf}). Is there a finite $k$ such that, for every $F\in\mathcal E$, the property that $f_F$ is a P\'olya frequency function of order $k$ implies that every zero of $F$ lies on the imaginary axis? \emph{Status.} Open. Chapter~\ref{chap:polyazeros} shows that the low orders are decided by the low zeros: order two holds without RH (Theorem~\ref{pol:thm:pf2}), and the density lies within $1.6\times10^{-10}$ of a P\'olya frequency function of every order (Theorem~\ref{pol:thm:approxpf}), so that, given the verification of RH to height $3\times10^{12}$~\cite{PT21}, every minor of order $3$ and $4$ is at least $-1.1\times10^{-8}$ and $-1.2\times10^{-7}$. Problem~\ref{pol:op:pf2} asks for proofs of the low orders, not for the converse asked here. Section~\ref{gl:sec:pfnum} reports the computations of this chapter. The perturbation of Theorem~\ref{thm:ch13:perturb}, made at a height that is large compared with $k$, is the natural candidate for a negative answer in the class $\mathcal E$; whether it keeps the order $k$ has not been computed here.
\end{problem}

\subsection*{Notes}

The Thorin representation and the van Dantzig pair are from~\cite{P17,P18}; the GGC background is Bondesson's monograph~\cite{Bondesson}; the master theorem class is from Chaudhry and Qadir~\cite{ChaudhryQadir12}, and the exponential Tauberian theorem from Kasahara~\cite{Kasahara78}. Completed Dirichlet $L$ functions and Dedekind zeta functions are in Davenport~\cite{Davenport}, mean values of $\zeta$ in Titchmarsh~\cite{Titchmarsh}. The account of the circle method in Section~\ref{gl:sec:where} is a summary of the paper of Heath-Brown and Puchta~\cite{HBP02}.
